\section{Применения суффиксного автомата и суффиксное дерево}

На прошлой лекции мы построили суффиксный автомат и доказали, что для строки длины $n$ он содержит не более $2n$ вершин и не более $3n$ рёбер. Сегодня мы завершим доказательство того, что алгоритм построения работает за линейное время $O(n)$, а затем увидим, как суффиксный автомат позволяет за $O(n)$ строить суффиксный массив и суффиксное дерево.

\subsection{Анализ времени построения суффиксного автомата}

Посмотрим на функцию \texttt{extend(c)}, добавляющую очередной символ $c$ в суффиксный автомат. Кроме операций, выполняющихся за $O(1)$, в ней есть три нетривиальных действия:
\begin{enumerate}
    \item \textbf{Первый цикл \texttt{while}}, проводящий переходы $p \xrightarrow{c} \text{cur}$, пока $\nxt(p, c) = -1$. На каждом шаге этого цикла в автомате появляется \textbf{новое ребро}. Поскольку рёбра из автомата никогда не удаляются (в дальнейшем они могут лишь перенаправляться в клон), суммарное число шагов первого цикла за все $n$ итераций в точности равно конечному числу рёбер $|E| \le 3n$.
    \item \textbf{Копирование исходящих переходов} при создании клона $\text{clone}$ ($\nxt(\text{clone}) = \nxt(q)$). Так как клон создаётся не более одного раза на символ, при фиксированном алфавите $\Sigma$ это занимает $O(n \cdot |\Sigma|)$ времени (а если хранить только существующие переходы в списке или дереве, каждое скопированное ребро увеличивает общее число рёбер $|E| \le 3n$, что даёт $O(n)$).
    \item \textbf{Второй цикл \texttt{while}}, перенаправляющий входящие переходы $p \xrightarrow{c} q$ в вершину $\text{clone}$. Здесь новые рёбра \emph{не создаются} "--- мы меняем конец уже существующих рёбер. Почему одно и то же ребро не может перенаправляться многократно, приводя к квадратичной работе?
\end{enumerate}

\begin{theorem}[О линейности числа перенаправлений рёбер]
    Суммарное количество итераций второго цикла \texttt{while} (перенаправлений переходов в вершину $\text{clone}$) за все $n$ шагов построения суффиксного автомата не превосходит $2n$.
\end{theorem}

\begin{proof}
    Для любой строки $s$ рассмотрим цепочку суффиксных ссылок, ведущую из состояния $[s]$ (отвечающего всей строке $s$) до корня $0$:
    \[
        L(s) = \bigl([s],\; \link([s]),\; \link(\link([s])),\; \dots,\; 0\bigr),
    \]
    и обозначим через $|L(s)|$ количество вершин в этой цепочке.

    Проследим, как связаны между собой цепочки $L(s)$ и $L(sc)$ при дописывании символа $c$ в Случае 3 (когда создаётся вершина $\text{clone}$ и выполняется второй цикл \texttt{while}).

    \begin{center}
    \begin{tikzpicture}[
        scale=1.0,
        vertex/.style={circle, draw=black!80, thick, minimum size=7.5mm, inner sep=2.5pt, font=\scriptsize},
        clonev/.style={vertex, fill=OliveGreen!20},
        newv/.style={vertex, fill=NavyBlue!20},
        linkedge/.style={draw=BrickRed, -Stealth, dashed, thick},
        trans/.style={draw=NavyBlue, -Stealth, thick}
    ]
        \node[font=\small\bfseries] at (-5.5, 1.5) {$L(s):$};
        \node[font=\small\bfseries] at (-5.5, -1.0) {$L(sc):$};

        % L(s) vertices
        \node[vertex] (s0) at (-4.0, 1.5) {$[s]$};
        \node[vertex] (s1) at (-2.4, 1.5) {$\dots$};
        \node[vertex] (p1) at (-0.6, 1.5) {$p$};
        \node[vertex] (p2) at (0.8, 1.5) {$\dots$};
        \node[vertex] (pk) at (2.2, 1.5) {$p'$};
        \node[vertex] (r1) at (3.8, 1.5) {$u_1$};
        \node[vertex] (r2) at (5.4, 1.5) {$0$};

        \draw[linkedge] (s0) -- (s1);
        \draw[linkedge] (s1) -- (p1);
        \draw[linkedge] (p1) -- (p2);
        \draw[linkedge] (p2) -- (pk);
        \draw[linkedge] (pk) -- (r1);
        \draw[linkedge] (r1) -- (r2);

        \draw[decorate, decoration={brace, amplitude=5pt}, thick] (-0.9, 2.0) -- (2.5, 2.0) node[midway, above=5pt, font=\footnotesize] {$k$ перенаправленных вершин};

        % L(sc) vertices
        \node[newv] (cur) at (-3.2, -1.0) {$[sc]$};
        \node[clonev] (clone) at (0.8, -1.0) {$\text{clone}$};
        \node[vertex] (w1) at (3.8, -1.0) {$w_1$};
        \node[vertex] (w2) at (5.4, -1.0) {$0$};

        \draw[linkedge] (cur) -- (clone);
        \draw[linkedge] (clone) -- (w1);
        \draw[linkedge] (w1) -- (w2);

        % Transitions by c
        \draw[trans] (s0) -- (cur);
        \draw[trans] (s1) -- (cur);
        \draw[trans, OliveGreen!80!black] (p1) -- (clone);
        \draw[trans, OliveGreen!80!black] (p2) -- (clone);
        \draw[trans, OliveGreen!80!black] (pk) -- (clone);
        \draw[trans] (r1) -- (w1);
    \end{tikzpicture}
    \end{center}

    Цепочка $L(sc)$ начинается с новой вершины $\text{cur} = [sc]$, за которой следует вершина $\text{clone} = \link(\text{cur})$, далее идут вершины $w_1, w_2, \dots$, и завершается цепочка корнем $0$. Любая вершина $w \in L(sc)$, кроме корня $0$, содержит непустые суффиксы строки $sc$, которые заканчиваются на букву $c$. Если отбросить у них последнюю букву $c$, мы получим суффиксы строки $s$, лежащие в вершинах цепочки $L(s)$.

    При этом из каждой вершины цепочки $L(s)$ выходит ровно один переход по букве $c$, ведущий в некоторую вершину цепочки $L(sc) \setminus \{0\}$:
    \begin{itemize}
        \item В первую вершину $[sc]$ ведёт хотя бы один переход из $L(s)$ (как минимум из вершины $[s]$).
        \item Во вторую вершину $\text{clone}$ ведут в точности те $k$ переходов из вершин $p, \link(p), \dots, p'$ цепочки $L(s)$, которые мы перенаправили во втором цикле \texttt{while}.
        \item В каждую из оставшихся $|L(sc)| - 3$ нетривиальных вершин цепочки $L(sc)$ (между $\text{clone}$ и корнем $0$) также входит хотя бы по одному переходу по букве $c$ из цепочки $L(s)$.
    \end{itemize}
    Так как все эти группы вершин лежат в $L(s)$ и попарно не пересекаются, общее число вершин в $L(s)$ не меньше суммы их размеров:
    \[
        |L(s)| \ge 1 + k + \bigl(|L(sc)| - 3\bigr) = |L(sc)| + k - 2.
    \]
    Перенося $k$ в левую часть, получаем оценку на изменение потенциала $|L|$:
    \[
        k \le |L(s)| - |L(sc)| + 2.
    \]
    Просуммируем это неравенство по всем $n$ шагам добавления символов (замечая, что в Случаях 1 и 2, когда $k = 0$, длина цепочки увеличивается не более чем на $1$: $|L(sc)| \le |L(s)| + 1$). Телескопическая сумма даёт
    \[
        \sum_{i=1}^n k_i \le |L(\epsilon)| - |L(s)| + 2n \le 1 - 1 + 2n = 2n.
    \]
    Тем самым доказано, что суффиксный автомат строится за линейное время $O(n)$.
\end{proof}

\subsection{Построение суффиксного массива за \texorpdfstring{$O(n)$}{O(n)}}

В качестве первого неожиданного применения суффиксного автомата покажем, как с его помощью построить суффиксный массив строки $s$ за линейное время $O(n)$ (улучшив оценку $O(n \log n)$ из третьей лекции).

Построим суффиксный автомат строки $s$ за $O(n)$ и пометим в нём все терминальные состояния (пройдя по суффиксным ссылкам от $[s]$ до корня $0$). Любой путь из корня $0$ в терминальную вершину отвечает некоторому суффиксу строки $s$, причём разным суффиксам отвечают разные пути.

Что будет, если мы запустим из корня $0$ обход в глубину (DFS), перебирающий исходящие из каждой вершины рёбра в алфавитном порядке букв, \textbf{не используя массив \texttt{visited}} (то есть заходя в каждую вершину столько раз, сколько различных путей ведёт в неё из корня)? Такой обход переберёт все пути из корня в лексикографическом порядке. Если при каждом заходе в терминальную вершину на глубине $d$ мы будем выписывать суффикс длины $d$ (начинающийся в позиции $n - d$), мы получим суффиксы строки $s$ в точности в порядке суффиксного массива.

Однако в исходном автомате такой обход без \texttt{visited} может работать за $\Theta(n^2)$: если в графе есть длинная цепочка вершин с одним исходящим ребром, в начало которой ведёт много путей, мы будем проходить всю эту цепочку заново для каждого пути.

Устраним эту проблему: \textbf{сожмём все нетерминальные проходные вершины}, из которых выходит ровно одно ребро. А именно, если нетерминальная вершина $u$ имеет единственный переход $u \xrightarrow{c} w$, мы можем заменить любое входящее в $u$ ребро $v \to u$ на прямое сжатое ребро $v \to w$, записав на нём:
\begin{itemize}
    \item \textbf{первую букву} этого сжатого пути (чтобы DFS знал, в каком порядке перебирать рёбра);
    \item \textbf{суммарную длину} $\Delta \ell$ сжатого пути в символах (чтобы при переходе по ребру увеличивать текущую длину прочитанной строки на $\Delta \ell$).
\end{itemize}
Сжатие всех таких вершин легко выполняется за один обход графа за $O(n)$. Оказывается, после этого сжатия обход в глубину без массива \texttt{visited} всегда работает за $O(n)$.

\begin{theorem}[О линейности обхода сжатого автомата]
    Обход в глубину (без пометок посещённости) по сжатому суффиксному автомату строки $s$ длины $n$ заходит во все вершины суммарно не более $3n + 1$ раз и строит суффиксный массив за время $O(n)$.
\end{theorem}

\begin{proof}
    Время работы обхода пропорционально суммарному числу заходов во все вершины автомата (с учётом кратности). Поскольку суффиксный автомат детерминирован, каждый заход в вершину $v$ однозначно задаётся строкой $\alpha$, прочитанной на пути от корня $0$ до $v$. Оценим отдельно число заходов в терминальные и нетерминальные вершины:

    \textbf{1. Заходы в терминальные вершины.} Каждый путь из корня в терминальную вершину задаёт суффикс строки $s$. Всего у строки $s$ ровно $n + 1$ суффиксов (включая пустой), поэтому в терминальные вершины обход зайдёт суммарно ровно $n + 1$ раз.

    \textbf{2. Заходы в нетерминальные вершины.} Пусть обход зашёл в некоторую нетерминальную вершину $v$ по пути со строкой $\alpha$. Так как мы сжали все нетерминальные вершины с исходящей степенью $1$ (а тупиковых нетерминальных вершин в суффиксном автомате не бывает), из вершины $v$ в сжатом автомате выходит \textbf{не менее двух рёбер}, начинающихся с различных символов $a \ne b$.

    Но это означает, что обе строки $\alpha a$ и $\alpha b$ являются подстроками строки $s$. Посмотрим теперь на \textbf{развёрнутую строку} $s^R$: в ней обе строки $a\alpha^R$ и $b\alpha^R$ являются подстроками, то есть подстрока $\alpha^R$ ветвится влево в строке $s^R$.

    По критерию самой длинной строки (лемма~\ref{lem:longest_crit}) отсюда следует, что строка $\alpha^R$ обязана быть самой длинной строкой ($\longest$) некоторого состояния в суффиксном автомате развёрнутой строки $s^R$. Но по теореме~\ref{thm:sa_size} в суффиксном автомате строки $s^R$ не более $2n$ вершин, а значит, существует не более $2n$ таких строк $\alpha^R$. Следовательно, во все нетерминальные вершины наш обход зайдёт суммарно не более $2n$ раз.

    Складывая две оценки, получаем не более $(n + 1) + 2n = 3n + 1$ заходов в вершины.
\end{proof}

\subsection{Суффиксное дерево и теорема о двойственности}

Наряду с суффиксным автоматом и суффиксным массивом третьей классической суффиксной структурой является \textbf{суффиксное дерево}.

\begin{definition}
    \textbf{Суффиксным деревом} строки $s$ называется сжатый бор, построенный по всем суффиксам строки $s$: в нём сжаты все нетерминальные вершины, из которых выходит ровно одно ребро, а на каждом сжатом ребре записана непустая подстрока строки $s$ (хранимая в виде пары индексов $[l, r]$ её начала и конца в строке $s$).
\end{definition}

Благодаря тому что на каждом ребре хранятся лишь два числа $[l, r]$, а не сами символы, размер каждого ребра в памяти составляет $O(1)$. Суффиксное дерево и суффиксный автомат тесно связаны друг с другом "--- фактически это две двойственные структуры.

\begin{theorem}[О двойственности суффиксного дерева и суффиксного автомата]\label{thm:st_sa_duality}
    Суффиксное дерево строки $s$ изоморфно дереву суффиксных ссылок суффиксного автомата развёрнутой строки $s^R$.
\end{theorem}

\begin{proof}
    Установим взаимно-однозначное соответствие между вершинами и рёбрами этих двух деревьев.

    \textbf{1. Биекция между вершинами.} Поймём, какие подстроки $\alpha$ строки $s$ остаются явными вершинами в боре всех суффиксов строки $s$ после сжатия однородных путей. Вершина $\alpha$ не сжимается тогда и только тогда, когда выполняется хотя бы одно из двух условий:
    \begin{itemize}
        \item вершина $\alpha$ является терминальной, то есть $\alpha$ "--- \textbf{суффикс} строки $s$;
        \item из вершины $\alpha$ в боре выходит не менее двух рёбер по разным буквам $a \ne b$, то есть обе строки $\alpha a$ и $\alpha b$ являются подстроками $s$ ($\alpha$ \textbf{ветвится вправо}).
    \end{itemize}
    Развернём теперь все строки задом наперёд, перейдя от строки $s$ к развёрнутой строке $s^R$ и от подстроки $\alpha$ к $\alpha^R$:
    \begin{itemize}
        \item условие <<$\alpha$ "--- суффикс $s$>> превращается в условие <<$\alpha^R$ "--- \textbf{префикс} $s^R$>>;
        \item условие <<$\alpha a$ и $\alpha b$ входят в $s$>> превращается в условие <<$a\alpha^R$ и $b\alpha^R$ входят в $s^R$>> ($\alpha^R$ \textbf{ветвится влево} в $s^R$).
    \end{itemize}
    Но по лемме~\ref{lem:longest_crit} это в точности критерий того, что строка $\alpha^R$ является самой длинной строкой ($\longest$) некоторого состояния суффиксного автомата строки $s^R$. Значит, отображение $\alpha \leftrightarrow \alpha^R$ задаёт биекцию между вершинами суффиксного дерева строки $s$ и состояниями суффиксного автомата строки $s^R$.

    \textbf{2. Биекция между рёбрами.} Пусть в суффиксном автомате строки $s^R$ суффиксная ссылка ведёт из состояния $u$ (с $\longest(u) = \alpha^R$) в состояние $v = \link(u)$ (с $\longest(v) = \beta^R$). По определению суффиксной ссылки, $\beta^R$ есть самый длинный собственный суффикс строки $\alpha^R$, который является строкой $\longest$ некоторого состояния в автомате для $s^R$.

    Переходя к развёрнутым строкам, получаем: $\beta$ есть самый длинный собственный \emph{префикс} строки $\alpha$, который является явной вершиной суффиксного дерева строки $s$. Но в сжатом боре самый длинный собственный префикс вершины $\alpha$, сам являющийся явной вершиной, "--- это в точности непосредственный родитель вершины $\alpha$ в суффиксном дереве. Значит, суффиксные ссылки в автомате для $s^R$ взаимно-однозначно соответствуют сжатым рёбрам суффиксного дерева строки $s$.
\end{proof}

\begin{corollary}
    Суффиксное дерево любой строки $s$ длины $n$ имеет не более $2n$ вершин и может быть построено за время $O(n)$ путём построения суффиксного автомата для развёрнутой строки $s^R$.
\end{corollary}

Чтобы при таком построении за $O(1)$ узнавать подстроку $s[l \dots r]$, написанную на ребре суффиксного дерева, достаточно в каждой вершине $u$ автомата для $s^R$ поддерживать позицию $\text{first\_pos}[u]$ конца первого появления класса $u$: для новой вершины $\text{cur}$ она равна текущей длине префикса, а при клонировании копируется из $q$ ($\text{first\_pos}[\text{clone}] = \text{first\_pos}[q]$).

\subsubsection{Сравнение суффиксного дерева и суффиксного автомата}

Хотя суффиксное дерево и суффиксный автомат двойственны друг другу, в конкретных задачах одна структура может быть заметно удобнее другой:
\begin{itemize}
    \item \textbf{Поиск $\LCP$ двух суффиксов за $O(1)$:} в суффиксном дереве у каждой вершины есть ровно один путь от корня, поэтому наибольший общий префикс двух суффиксов $s^i$ и $s^j$ "--- это в точности наименьший общий предок ($\LCA$) двух соответствующих терминальных вершин дерева. С предподсчётом за $O(n)$ (сведением $\LCA$ к $\RMQ$ на эйлеровом обходе) запросы $\LCP$ выполняются за $O(1)$. В суффиксном автомате в одну вершину ведут пути разной длины, поэтому $\LCA$ к нему напрямую неприменим.
    \item \textbf{Проверка совпадения множеств вхождений двух подстрок $t_1$ и $t_2$:} в суффиксном автомате достаточно прочитать обе строки из корня и проверить, привели ли они в одно и то же состояние (и совпадают ли их правые контексты).
    \item \textbf{Поддержание скользящего окна по тексту:} пусть мы двигаем окно $t[l \dots r]$ по тексту $t$, умея при сдвиге $r \to r+1$ дописывать символ справа, а при сдвиге $l \to l+1$ "--- отбрасывать первый символ слева, и хотим знать, входит ли текущее окно в строку $s$. В суффиксном автомате удаление первого символа у текущей подстроки длины $L$, находящейся в вершине $v$, делается тривиально: мы уменьшаем текущую длину $L \la L - 1$, и если она стала равна $\len(\link(v))$, переходим по суффиксной ссылке $v \la \link(v)$, а иначе просто остаёмся в той же вершине $v$. В суффиксном дереве мы бы находились посреди сжатого ребра, и переход требовал бы дополнительных прыжков.
\end{itemize}

\begin{note}[Суффиксные ссылки в суффиксном дереве и жёсткие рёбра]
    В самом суффиксном дереве тоже определяют суффиксные ссылки: для явной вершины $\alpha = c\gamma$ её суффиксная ссылка ведёт в строку $\gamma$, полученную отбрасыванием первой буквы $c$. Легко видеть, что $\gamma$ всегда является \emph{явной} вершиной суффиксного дерева (если $c\gamma$ была суффиксом $s$, то и $\gamma$ "--- суффикс $s$; если $c\gamma a$ и $c\gamma b$ входили в $s$, то и $\gamma a, \gamma b$ входят в $s$).

    В силу теоремы~\ref{thm:st_sa_duality} о двойственности переход $\gamma \to c\gamma$ в суффиксном дереве строки $s$ соответствует в суффиксном автомате развёрнутой строки $s^R$ переходу $\gamma^R \xrightarrow{c} \gamma^R c = \alpha^R$, при котором длина самой длинной строки класса увеличивается ровно на единицу:
    \[
        \len(\alpha^R) = \len(\gamma^R) + 1.
    \]
    Такие переходы $p \xrightarrow{c} q$ в суффиксном автомате, для которых $\len(p) + 1 = \len(q)$, называются \textbf{жёсткими рёбрами}. По лемме~\ref{lem:sa_incoming} в каждую нетривиальную вершину суффиксного автомата входит ровно одно жёсткое ребро "--- и при развороте оно в точности задаёт суффиксную ссылку в суффиксном дереве.
\end{note}

\subsection{Обобщённый суффиксный автомат для набора строк}

В биоинформатике и анализе данных часто возникает задача одновременной работы с набором из нескольких строк $s_1, s_2, \dots, s_k$ (например, поиск наибольшей общей подстроки нескольких геномных последовательностей ДНК).

Суффиксный автомат естественным образом обобщается на произвольный набор строк $s_1, \dots, s_k$. Для этого достаточно определить правый контекст строки $\alpha$ как множество пар:
\[
    R(\alpha) = \bigl\{(\beta, i) \;\big|\; 1 \le i \le k,\ \alpha\beta \text{ является суффиксом строки } s_i\bigr\}.
\]
Все структурные леммы (о вложенности контекстов, интервале суффиксов в каждом классе и критерии $\longest$) дословно сохраняют силу.

Чтобы построить \textbf{обобщённый суффиксный автомат} по строкам $s_1, \dots, s_k$, достаточно:
\begin{enumerate}
    \item построить обычный суффиксный автомат для первой строки $s_1$;
    \item перед добавлением каждой следующей строки $s_i$ ($i \ge 2$) обнулять текущую вершину $\text{last} = 0$ (возвращаясь в корень) и посимвольно вызывать функцию \texttt{extend} для символов строки $s_i$.
\end{enumerate}
Единственная техническая деталь при таком добавлении состоит в том, что из текущей вершины $\text{last}$ переход по очередному символу $c$ может \emph{уже существовать} в автомате (если префикс новой строки совпал с подстрокой предыдущих строк). В этом случае создавать новую вершину $\text{cur}$ не нужно: если для существующего перехода $\text{last} \xrightarrow{c} q$ выполнено $\len(\text{last}) + 1 = \len(q)$, мы просто переходим $\text{last} \la q$; если же $\len(\text{last}) + 1 < \len(q)$, мы отщепляем от $q$ вершину $\text{clone}$ длины $\len(\text{last}) + 1$ и переходим $\text{last} \la \text{clone}$. Суммарное время построения обобщённого автомата линейно по сумме длин всех строк: $O\!\left(\sum_{i=1}^k |s_i|\right)$.
